\chapter{Desarrollo experimental y de ingeniería}\label{chap:desarrollo}
Este capítulo expone la arquitectura lógica, el flujo de orquestación y las garantías de reproducibilidad del sistema, de forma que cualquier resultado—figura, tabla o métrica—pueda regenerarse a partir de una configuración cerrada, artefactos versionados y un historial auditable de decisiones \cite{Tsamardinos2022,Vabalas2019,EAFIT2018}.

\section{Arquitectura funcional y flujo}
El sistema se organiza como un \textit{pipeline} dirigido por configuración con contratos explícitos entre etapas y artefactos inmutables. La ejecución sigue un recorrido determinista: configuración efectiva \(\rightarrow\) verificación y normalización \(\rightarrow\) representación de hologramas \(\rightarrow\) aprendizaje y calibración \(\rightarrow\) validación y reporte \(\rightarrow\) seguridad de dominio. La \emph{configuración efectiva} se consolida desde un archivo declarativo (\texttt{config.yaml}) y fija semillas de NumPy y \texttt{scikit-learn}, además de congelar versiones de dependencias para instaurar determinismo a nivel de librerías \cite{Pedregosa2011}. La etapa de \emph{verificación y normalización} aplica las operaciones radiométricas y geométricas definidas en el Capítulo~\ref{chap:metodologia}—percentiles \(P_{1}\)–\(P_{99}\), tamaño objetivo \(512\times512\), CLAHE, filtrado de mediana y compensación de tilts/curvaturas en frecuencia—y emite un manifiesto con estadísticas pre/post y huellas criptográficas que anclan el linaje de los datos \cite{Greenbaum2014,Amann2019,OConnor2020,Singh2017}. La \emph{representación de hologramas} construye vectores \(\mathbf{x}\) a partir de LBP y GLCM, energía en anillos de Fourier y anisotropía espectral, bancos de Gabor, energías wavelet y morfometría; todos los descriptores se estandarizan con estimadores robustos y se someten a reducción univariada confinada al entrenamiento interno para evitar fuga de información \cite{Tsamardinos2022,Vabalas2019}.

El bloque de \emph{aprendizaje y calibración} ajusta clasificadores complementarios (regresión logística penalizada, SVM RBF, bosque aleatorio y potenciación de gradiente) y combina sus probabilidades en una mezcla convexa cuyos pesos se aprenden para minimizar la pérdida de Brier; la calibración posterior (isotónica o logística) se realiza en un pliegue independiente del usado para la búsqueda de hiperparámetros, preservando la independencia entre etapas \cite{Pedregosa2011,Friedman2001}. La \emph{validación y reporte} estima ROC-AUC, PR-AUC cuando procede, precisión balanceada, sensibilidad y especificidad con intervalos de confianza por \emph{bootstrap} estratificado, y genera curvas ROC/PR y diagramas de confiabilidad bajo un esquema anidado que confina selección y escalado al entrenamiento \cite{Vabalas2019,Tsamardinos2022}. Finalmente, la \emph{seguridad de dominio} entrena un detector de Mahalanobis sobre el subespacio activo, fija umbrales frente a \(\chi^{2}_{d}\) y calcula contribuciones por característica para explicar alertas; su evaluación usa hologramas externos no sanguíneos, cuantificando tasas de verdadera y falsa alarma y validando la estabilidad de umbrales por remuestreo \cite{Mahalanobis1936,Rousseeuw2019}.

\section{Automatización, trazabilidad y aseguramiento de calidad}
La orquestación se expone mediante una interfaz de línea de comandos (\texttt{cli.py}) con semántica centrada en tareas. El subcomando \texttt{analyze} ejecuta el flujo de punta a punta con la \emph{configuración efectiva} (rutas, parámetros de normalización, radios/orientaciones de descriptores, \(k\) de selección, hiperparámetros y semillas); \texttt{validate} procesa directorios externos para evaluación y anomalías; \texttt{info} y \texttt{config} entregan diagnósticos (versiones, semillas, tamaños y rutas efectivas) útiles para auditoría rápida. Cada corrida produce un paquete de resultados con marca temporal que incluye: volcado íntegro de la configuración utilizada, artefactos serializados (vectorizadores, modelos, calibradores y detector) etiquetados con hashes de contenido, bitácoras de ejecución (tiempos, recursos y semillas) y productos analíticos (figuras y tablas). Las etapas intensivas—unificación, extracción y selección—son idempotentes: si la huella \(\mathrm{hash}(\text{datos}+\text{config})\) coincide, se reutiliza el artefacto; si difiere, se regenera y se registra la nueva huella. Este diseño reduce cómputo redundante y hace explícito cuándo y por qué cambian los resultados.

El aseguramiento de calidad combina pruebas de humo que recorren la cadena completa sobre un subconjunto mínimo con verificaciones de contratos (selección confinada al pliegue, fijación de semillas antes de cualquier partición, consistencia de tamaños tras recorte y normalización). El control de versiones registra cambios atómicos y fija \texttt{requirements.txt} para preservar determinismo entre instalaciones \cite{Pedregosa2011,EAFIT2018}. En conjunto, la arquitectura—dirigida por configuración, con estados idempotentes, separación estricta entre entrenamiento, búsqueda, calibración y prueba, y evidencia versionada—materializa las garantías de reproducibilidad, trazabilidad y gobernanza requeridas en investigación aplicada en salud, sin depender de detalles contingentes de la disposición física de carpetas \cite{Tsamardinos2022,Vabalas2019,EAFIT2018}.
